\documentclass[12pt]{article}

% ------------------------------------------------------------
% Packages
% ------------------------------------------------------------
\usepackage[margin=1in]{geometry}
\usepackage{mathpazo}
\usepackage{amsmath,amssymb,amsfonts}
\usepackage{nicefrac}
\usepackage{microtype}
\usepackage{hyperref}

\hypersetup{
    colorlinks=true,
    linkcolor=blue,
    urlcolor=blue
}

% ------------------------------------------------------------
% Custom commands
% ------------------------------------------------------------
\newcommand{\N}{\mathbb{N}}
\newcommand{\Z}{\mathbb{Z}}
\newcommand{\Q}{\mathbb{Q}}
\newcommand{\R}{\mathbb{R}}

% ------------------------------------------------------------
% Title
% ------------------------------------------------------------
\title{Math 714 — Assignment 0}
\author{Tyler Jones}
\date{Fall 2026}

% ------------------------------------------------------------
\begin{document}
\maketitle

\section*{Instructions}
This document contains my solutions and computational work for Assignment~0 of Math~714.  
All problem statements are reproduced below for clarity, followed by my solutions.

\bigskip

% ============================================================
% Problem 1
% ============================================================
\section*{Problem 1}
Consider the recursive sequence


\[
x_{k+1} = a x_k^2 + b x_k, \qquad k = 0,1,2,\ldots
\]



\begin{enumerate}
    \item[(a)] If the sequence converges, to what value or values does it converge?
    \item[(b)] Suppose that $x_0$ is arbitrarily close to the value or values above.  
    For what parameters $(a,b)$ does the sequence $x_k$ converge to that value?
\end{enumerate}

\subsection*{Solution}
% --- Your solution goes here ---


% ============================================================
% Problem 2
% ============================================================
\section*{Problem 2}
The Chebyshev polynomials $T_k(x)$ satisfy the recurrence


\[
T_k(x) = 2x\,T_{k-1}(x) - T_{k-2}(x),
\qquad T_0(x)=1,\; T_1(x)=x.
\]



Evaluate and plot $T_5(x)$ at 101 evenly spaced points in $[-1,1]$.  
Draw a 2D surface plot of $T_3(x)T_5(y)$ on a $101\times 101$ grid over $(x,y)\in[-1,1]^2$.

\subsection*{Solution}
% --- Your solution goes here ---


% ============================================================
% Problem 3
% ============================================================
\section*{Problem 3}
The IEEE double-precision standard guarantees that for any operation $\ast$,  
the floating point operation $\circledast$ satisfies


\[
x \circledast y = (1+\delta)(x \ast y), \qquad |\delta| < \epsilon.
\]



\begin{enumerate}
    \item[(a)] Compute the minimum and maximum possible values of
    

\[
    S = \frac{3}{4+2}
    \]


    when evaluated using floating point arithmetic as
    

\[
    \tilde{S} = 3 \oslash (4 \oplus 2).
    \]



    \item[(b)] Show that if $O(\epsilon^2)$ terms are neglected, then
    

\[
    |S - \tilde{S}| < \lambda \epsilon,
    \]


    and determine the constant $\lambda$.
\end{enumerate}

\subsection*{Solution}
% --- Your solution goes here ---


% ============================================================
% Problem 4
% ============================================================
\section*{Problem 4}
For an invertible matrix $A$, define the condition number


\[
\kappa(A) = \|A\|\,\|A^{-1}\|.
\]



Assume the matrix norm is induced by the Euclidean vector norm.

\begin{enumerate}
    \item[(a)] Find $2\times 2$ invertible matrices $B$ and $C$ such that  
    $\kappa(B+C) < \kappa(B) + \kappa(C)$.

    \item[(b)] Find $2\times 2$ invertible matrices $B$ and $C$ such that  
    $\kappa(B+C) > \kappa(B) + \kappa(C)$.

    \item[(c)] Suppose $A$ is symmetric and invertible.  
    Find $\kappa(A^2)$ in terms of $\kappa(A)$.

    \item[(d)] Does the result from part (c) hold if $A$ is not symmetric?  
    Either prove the result or give a counterexample.

    \item[(e)] For invertible matrices $B$ and $C$, prove that  
    $\kappa(BC) \le \kappa(B)\kappa(C)$.  
    Give examples where equality holds and where the inequality is strict.
\end{enumerate}

\subsection*{Solution}
% --- Your solution goes here ---


% ============================================================
% Problem 5
% ============================================================
\section*{Problem 5}
The gamma function is defined by


\[
\Gamma(x) = \int_0^\infty t^{x-1} e^{-t}\,dt.
\]



It satisfies $(n-1)! = \Gamma(n)$ for integers $n$.  
Also,


\[
\Gamma\!\left(\tfrac32\right) = \tfrac12\sqrt{\pi}.
\]



\begin{enumerate}
    \item[(a)] For $k=1,\ldots,5$, compute polynomials $p_k(x)$ of degree $k$  
    matching $\Gamma(x)$ at $x=1,2,\ldots,k+1$.  
    Evaluate the error $\bigl|p_k(\tfrac32) - \Gamma(\tfrac32)\bigr|$.

    \item[(b)] For $k=1,\ldots,5$, compute polynomials $q_k(x)$ of degree $k$  
    matching $\log\Gamma(x)$ at $x=1,2,\ldots,k+1$.  
    Evaluate the error $\bigl|\exp(q_k(\tfrac32)) - \Gamma(\tfrac32)\bigr|$.

    \item[(c)] Examine the asymptotic behavior of both errors as $k\to\infty$.
\end{enumerate}

\subsection*{Solution}
% --- Your solution goes here ---


% ============================================================
% Problem 6
% ============================================================
\section*{Problem 6}

\begin{enumerate}
    \item[(a)] Let $f:\R\to\R$ be five-times differentiable.  
    Expand Taylor series for $f(x-h)$ and $f(x+h)$ to derive coefficients  
    $\alpha,\beta,\gamma$ such that
    

\[
    f''(x) = \frac{\alpha f(x-h) + \beta f(x) + \gamma f(x+h)}{h^2} + O(h).
    \]



    \item[(b)] Suppose $f'$ can be evaluated exactly.  
    Find coefficients $a,b,c,r,s,t$ such that
    \begin{align*}
        f''(x)
        &= \frac{a f(x-h) + b f(x) + c f(x+h)}{h^2} \\
        &\quad + \frac{r f'(x-h) + s f'(x) + t f'(x+h)}{h}
        + O(h^4).
    \end{align*}

    \item[(c)] Consider $f(x) = e^{4\sin x}$.  
    Compute $f'$ and $f''$ analytically.  
    Write a program to test the formulas in (a) and (b) at $x=1$  
    using $h = 2^{-k}$ for $k=0,1,\ldots,23$.  
    Make log–log plots of the error $E$ vs.\ $h$ and fit $E = Ch^p$.  
    Discuss whether the observed $p$ matches the theoretical order.
\end{enumerate}

\subsection*{Solution}
% --- Your solution goes here ---


% ------------------------------------------------------------
\end{document}
